\begin{ExerciseBlock}
\ExerciseHeading{03}{Schiefer Wurf: Optimaler Abwurfwinkel}{multidimensional_kinematics_optimal_launch_angle}
Ein Weitwerfer wirft einen Ball vom Boden aus im Winkel $\alpha$ mit der Geschwindigkeit $v_0$. Luftwiderstand wird vernachlässigt.

\begin{enumerate}
	\item Stellen Sie die Ortsfunktion $\vec{r}(t)$ für den Ball auf. (Beachten Sie außerdem, dass die Geschwindigkeit in Vektorkomponenten
zerlegt werden muss.)

	\item Leiten Sie folgende Formel für die Zeitdauer bis zum Aufschlag auf den Boden her: \[ t_\text{Wurf} = \frac{2 v_0 \sin(\alpha)}{g} \]

	\item Zeigen Sie, dass für die Wurfdistanz $s_\text{Wurf}$ gilt: \[ s_\text{Wurf} = \frac{v_0^2}{g} \sin(2\alpha) \]
	Hinweis: Für den Sinus gilt das Additionstheorem \[ \sin(\alpha+\beta) = \sin(\alpha) \cos(\beta) + \cos(\alpha) \sin(\beta) \]
	Für $\alpha = \beta$ folgt $\sin(2 \alpha) = 2 \sin(\alpha) \cos(\alpha)$.

	\item Bestimmen Sie den optimalen Abwurfwinkel $\alpha_{opt}$, um maximale Wurfweite zu erzielen.

\end{enumerate}
\end{ExerciseBlock}
\UseExerciseSolution{multidimensional_kinematics_optimal_launch_angle}
